% !TeX program = xelatex
\input{preamble.tex}
\newcommand{\RR}{\mathcal R}
\newcommand{\FF}{\mathcal F}
\newcommand{\GG}{\mathcal G}
\newcommand{\MM}{\mathcal M}
\newcommand{\JJ}{\mathcal J}
\newcommand{\oldauditsrc}[1]{\par\noindent\textbf{原文位置：}#1。\par}
\newcommand{\kind}[2]{\noindent\textbf{分类：#1。}\ #2\par}
\newcommand{\lead}[1]{\par\noindent\textbf{#1}\quad}
\newenvironment{analysisproof}{\par\noindent\textbf{该轮核查的反例或补证。}\ }{\par}

\hypersetup{pdftitle={代数几何基础：第十一版译校核查报告},pdfsubject={第0章21.1--21.4续译定点核查；原作者勘误；历史核查档案}}
\title{代数几何基础\\译校核查报告\\[1em]\large 第十一版新核查与历史记录}
\author{原著对照、反例与补证}
\date{}
\begin{document}
\frontmatter
\maketitle
\head{报告说明}
本报告分为本轮新核查与历史记录。本轮依据上传法文原页，接续第0章§21章首及21.1--21.4，止于（21.4.8）；新标识以K11开头。
新增正文写入经核实的校订表述，原文差异、修改依据置于脚注；详细反例与补证收在本报告。
重点包括$p$-无关的系数环、集合并与有索引的族、Frobenius半线性、相对微分模的基环，
以及$p$-基不自动保证形式光滑等使用边界。

历史报告保留当时编号和措辞，不表示本轮重新审定其所有结论。
旧版H8-P01等未决事项仍按原状态保留。
合订本仍是当前已译内容的工作稿，不是全部上传材料的完结全译。
本报告中的新反例与修订均单独给出论证，不冒称为作者正式勘误。
编译、条目核对和渲染检查不代替出版级逐句审校或全书数学证明的独立验证。
\tableofcontents
\mainmatter
\input{reports/eleventh_audit.tex}
\part*{历史核查记录：第十版}
\addcontentsline{toc}{part}{历史核查记录：第十版}
\NewNote{以下J10条目保留第十版的核查判断及范围说明，不表示本轮重新审定其全部内容。}
\input{reports/tenth_audit.tex}
\part*{历史核查记录：第九版}
\addcontentsline{toc}{part}{历史核查记录：第九版}
\NewNote{以下I9条目保留第九版的定点判断和范围说明，不表示本轮已重新审定全部历史核查。}
\input{reports/ninth_audit.tex}
\part*{历史核查记录：第八版}
\addcontentsline{toc}{part}{历史核查记录：第八版}
\NewNote{以下H8条目保留第八版原有判断和范围说明。本轮未重新审定全部旧稿；H8-P01的拓扑范围仍属待核事项。}
\input{reports/eighth_audit.tex}
\part*{历史核查记录：第七版}
\addcontentsline{toc}{part}{历史核查记录：第七版}
\NewNote{以下G7条目为第七版历史记录。其内容保留，不表示本轮重新逐项审定；本版新的拓扑范围说明以H8-P01、H8-P02为准。}
\input{reports/seventh_audit.tex}
\part*{历史核查记录：第六版}
\addcontentsline{toc}{part}{历史核查记录：第六版}
\NewNote{以下F6条目为第六版历史记录。其中关于正文仍保留错误原证明的描述已被本版相应修订取代；本版并未重新审定其余所有历史判断。}
\input{reports/sixth_audit.tex}
\part*{历史核查记录：第五版}
\addcontentsline{toc}{part}{历史核查记录：第五版}
\NewNote{以下E5条目为第五版原有记录，不表示本轮重新审定其全部结论。}
\input{reports/fifth_audit.tex}
\part*{历史核查记录：第四版}
\addcontentsline{toc}{part}{历史核查记录：第四版}
\NewNote{以下D4条目为第四版原有记录，不表示本轮重新审定其全部结论。}
\input{reports/fourth_audit.tex}
\part*{历史核查记录：第三版}
\addcontentsline{toc}{part}{历史核查记录：第三版}
\NewNote{下列C3条目属于第三版的定点核查，不是本轮重新发现或再次逐项独立审定的结果。}
\input{reports/third_order_audit.tex}
\part*{历史核查记录：前次续译批次}
\addcontentsline{toc}{part}{历史核查记录：前次续译批次}
\NewNote{以下为上一批报告的原有内容；其“本轮”指当时批次，不是本第十一版新核查。}
\input{reports/current_audit.tex}
\part*{历史核查记录：新增文稿第一批}
\addcontentsline{toc}{part}{历史核查记录：新增文稿第一批}
\NewNote{以下为新增文稿第一批报告，保留原标识及其适用范围。}
\input{reports/batch1_audit.tex}
\part*{历史核查记录：第一、第二卷}
\addcontentsline{toc}{part}{历史核查记录：第一、第二卷}
\NewNote{以下为第一、第二卷早期重点审阅结果。关于进制完备性和对角态射等问题，须结合本版新取得的ErrIII清单；旧报告不代表问题尚未有正式勘误。}
\input{reports/legacy_I_II.tex}
\end{document}
